\chapter{Verified Embeddings and Stratified Soundness}
\label{ch:39}

In Chapters~31 and~32 we constructed the semantic embedding
[
\iota : \mathrm{Inc}(\mathcal{G}) \hookrightarrow \mathcal{M},
]
defined for any Polyxan hypergraph $\mathcal{G}$ as the unique harmonic, curvature-matched map that satisfies the embedding energy minimisation principle.

We now prove that this embedding is verified in the precise sense of Part~IV: it satisfies the two modal correctness predicates
[
\Box_{\mathrm{RSVP}} \qquad\text{and}\qquad \Box_{\mathrm{Polyx}}.
]

This chapter establishes \emph{stratified soundness}—the guarantee that the embedding behaves correctly on every curvature stratum of $\mathcal{M}$, and that its evolution under the embedding flow preserves harmonicity, quine-stability, and total correctness.

\section{Stratified Embeddings}

The embedding $\iota$ is stratified because both the source (the incidence complex) and the target (the semantic manifold) possess natural stratifications.

\begin{definition}[Stratified Embedding]
An embedding
[
\iota : \mathrm{Inc}(\mathcal{G}) \to \mathcal{M}
]
is \emph{stratified} if:

\begin{enumerate}
\item Each cell $c\in \mathrm{Inc}(\mathcal{G})$ is mapped into a unique curvature stratum
[
\mathcal{M}_k = {x\in\mathcal{M} : \mathcal{R}(x)=k}.
]

\item The lift of the primitive curvature satisfies
[
\mathcal{R}\circ\iota = \mathcal{H}_0[\mathcal{G}] \qquad\text{on each cell}.
]

\item The local entropy satisfies
[
S\circ\iota = \mathcal{H}_{\mathrm{tag}}[\mathcal{G}].
]
\end{enumerate}
\end{definition}

By the curvature–entropy matching theorem (Ch.~31), the embedding constructed by the RSVP flow is the unique stratified embedding.

\section{Harmonicity and Modal Necessity}

We now link the analytic property of being harmonic to the modal operator $\Box_{\mathrm{RSVP}}$.

\begin{lemma}[Harmonic Necessity]
If $\iota$ is harmonic, then
[
\Box_{\mathrm{RSVP}}(\iota),
]
that is, $\iota$ is provably preserved under all harmonic evolutions.
\end{lemma}

\begin{proof}
By definition,
[
\Box_{\mathrm{RSVP}} P
:= \forall \phi,; \text{HarmonicEvolution}(\phi) \to P(\phi).
]
A harmonic embedding $\iota$ is a fixed point of the Lyapunov energy functional (Ch.~32).
Every harmonic evolution preserves all fixed points.
Thus for any evolution beginning at $\iota$, the property “is harmonic and curvature-matched’’ remains true, which is exactly the predicate guarded by the modal operator.
\end{proof}

\section{Quine-Stability and Modal Necessity}

Similarly, the Polyxan modality captures the invariants of EHR rewriting.

\begin{lemma}[Quine-Stability Necessity]
If $\mathcal{G}$ is quine-stable, then
[
\Box_{\mathrm{Polyx}}(\iota).
]
\end{lemma}

\begin{proof}
A quine-stable graph admits no further EHR steps.
The embedding $\iota$ inherits this fixedness under the synthetic duality (Ch.~62).
By definition, $\Box_{\mathrm{Polyx}}$ asserts invariance under all EHR steps, which holds trivially when the underlying graph is already fixed.
\end{proof}

\section{Stratified Soundness Theorem}

We now assemble the pieces into the main theorem of the chapter.

\begin{theorem}[Stratified Soundness]
\label{thm:39-soundness}
Let $\mathcal{G}$ be any Polyxan hypergraph and let
[
\iota : \mathrm{Inc}(\mathcal{G}) \hookrightarrow \mathcal{M}
]
be the harmonic, curvature-matched embedding produced by the embedding flow.
Then
[
\Box_{\mathrm{RSVP}}(\iota)
\qquad\text{and}\qquad
\Box_{\mathrm{Polyx}}(\iota).
]
In other words:
[
\iota \text{ is verified.}
]
\end{theorem}

\begin{proof}
There are three steps.

\medskip
\noindent\textbf{1. Harmonicity.}
By construction (Ch.~32), $\iota$ is a minimizer of the embedding energy functional.
Thus it satisfies the harmonic map equation.
By the Harmonic Necessity Lemma, this implies $\Box_{\mathrm{RSVP}}(\iota)$.

\medskip
\noindent\textbf{2. Curvature–Entropy Matching.}
Because
[
\mathcal{R}\circ\iota = \mathcal{H}*0[\mathcal{G}]
\quad\text{and}\quad
S\circ\iota = \mathcal{H}*{\mathrm{tag}}[\mathcal{G}],
]
every EHR rewrite (if any) is immediately mirrored by a topology-changing event in the embedding flow (Ch.~33).
Thus $\iota$ evolves in modal lockstep with $\mathcal{G}$.

\medskip
\noindent\textbf{3. Quine-Stability.}
When $\mathcal{G}$ converges to its normal form, the embedding flow simultaneously converges to a harmonic representative supported on a flat curvature stratum.
By the Quine-Stability Necessity Lemma, this proves $\Box_{\mathrm{Polyx}}(\iota)$.

\medskip
Combining all three steps yields the result.
\end{proof}

\section{Consequences for the Coupled System}

The Stratified Soundness Theorem establishes the full correctness of the embedding layer.
Important consequences include:

\begin{corollary}[Inviolability of Harmonicity]
No verified dynamic step (Ch.~37) can break harmonicity.
Any admissible evolution preserves the curvature–entropy matching of $\iota$.
\end{corollary}

\begin{corollary}[Inviolability of Quine-Stability]
If $\iota$ is quine-stable, no verified Polyxan rewrite can move it out of its stability basin.
\end{corollary}

\begin{corollary}[End-to-End Verification]
For any coupled configuration,
[
(\Phi,\mathbf{v},S;,\mathcal{G},\iota),
]
the verified evolution preserves both modalities and thus preserves correctness of the embedding layer at all times.
\end{corollary}

\section{Closing Cadence}

The embedding is more than a geometric device.
It is the object that unites the continuous RSVP fields and the discrete Polyxan syntax.
By proving its stratified soundness under both modalities, we obtain:

\begin{quote}
The embedding cannot err.
It evolves only along verified trajectories,
and it encodes meaning in a form that cannot be broken.
\end{quote}

